\section*{Глава 1. Анализ предметной области}
\addcontentsline{toc}{section}{Глава 1. Анализ предметной области}
\stepcounter{section}

\subsection{Обзор существующих приложений для учета питания}

На современном рынке цифровых решений для контроля питания можно выделить две основные группы: мобильные приложения с облачной синхронизацией и веб-ориентированные платформы для анализа рациона. Проведем сравнение наиболее популярных решений в таблице ~\ref{tab:competitors}.

\begin{table}[ht!]
	\centering
	\caption{Сравнительный анализ приложений для учета питания}
	\label{tab:competitors}
	{\small
\begin{tabular}{|l|l|l|l|}
	\hline
	\textbf{Критерий} & \textbf{MyFitnessPal} & \textbf{Yazio} & \textbf{FatSecret} \\ \hline
	Ключевая особенность & Обширная база продуктов & Планы питания & Соц. функции \\ \hline
	Модель доступа & Freemium & Freemium & Бесплатно \\ \hline
	Наличие рекламы & Есть & Есть & Нет \\ \hline
	Готовые планы питания & Только премиум & Премиум & Нет \\ \hline
	Графики и аналитика & Базовая & Детальная & Базовая \\ \hline
	Русскоязычный интерфейс & Нет & Да & Да \\ \hline
	Социальные функции & Да & Нет & Да \\ \hline
\end{tabular}
}
\end{table}

\subsection*{Общие недостатки существующих решений:}
\begin{itemize}
	\item Большинство приложений используют модель Freemium — расширенные функции (детальная аналитика, планы питания) доступны только за плату;
	\item Наличие рекламы в бесплатных версиях, что ухудшает пользовательский опыт;
	\item Отсутствие гибкости настройки под индивидуальные потребности пользователя;
	\item Ограниченные возможности по управлению и расширению справочника продуктов через веб-интерфейс.
\end{itemize}

\textbf{Актуальность разработки собственного решения:}

Таким образом, создание веб-приложения для подсчета калорий и анализа питания, которое будет полностью бесплатным, не содержать рекламы, предоставлять возможность управления справочником продуктов через административную панель и обеспечивать базовую аналитику по калорийности и БЖУ, является актуальной задачей. Разрабатываемое приложение также позволит пользователю вести дневник питания без ограничений и адаптировать систему под свои индивидуальные потребности.

\subsection{Анализ программных инструментов разработки}

Для реализации проекта рассматриваются современные подходы к разработке веб-приложений с клиент-серверной архитектурой. Выбор инструментов обусловлен требованиями к удобству интерфейса, производительности и масштабируемости системы:

\begin{itemize}
	\item \textbf{Frontend-фреймворк:} Использование тяжелых фронтенд-фреймворков (React, Vue.js) и подход на основе чистого JavaScript с Bootstrap. React и Vue.js являются мощными инструментами, предоставляющими компонентную архитектуру, виртуальный DOM и богатые экосистемы. Однако их применение в данном проекте признано избыточным по следующим причинам: интерфейс дневника питания представляет собой совокупность статичных страниц с ограниченной динамикой; использование React или Vue.js потребовало бы настройки системы сборки (Webpack, Vite), изучения дополнительных концепций (хуки, маршрутизация) и увеличило бы время разработки без существенного выигрыша в функциональности. В связи с этим выбран подход на основе HTML, CSS, Bootstrap и чистого JavaScript.
	\item \textbf{Backend-платформа:} Для реализации серверной части был проведен сравнительный анализ FastAPI, Flask и Django. FastAPI является современным асинхронным фреймворком с высокой производительностью и автоматической генерацией документации, однако он не предоставляет встроенной ORM, административной панели и системы аутентификации, что потребовало бы написания значительного объема дополнительного кода. Flask — легковесный микрофреймворк, обладающий гибкостью и большим количеством расширений, но также требующий самостоятельной реализации многих компонентов. Django, в отличие от них, предлагает встроенную ORM для удобной работы с PostgreSQL, готовую административную панель для управления справочником продуктов, встроенную систему аутентификации пользователей и механизм миграций для автоматического обновления структуры базы данных. Эти преимущества существенно ускоряют разработку и делают Django оптимальным выбором для данного проекта. Использование встроенной ORM позволяет эффективно обрабатывать запросы от множества пользователей при одновременном подсчете калорий и агрегации данных о питании за длительные периоды. Bootstrap обеспечивает быструю адаптивную верстку и готовые компоненты интерфейса (формы, таблицы, карточки), а чистый JavaScript позволяет реализовать необходимую динамику: отправку асинхронных запросов к серверу, динамическое обновление таблицы приемов пищи и валидацию вводимых данных без избыточного усложнения проекта.
    \item \textbf{Хранение данных:} Для хранения данных рассматривались реляционная СУБД PostgreSQL и легковесная SQLite. SQLite не требует установки и настройки, что удобно для прототипирования, однако имеет ограниченную поддержку сложных запросов, снижает производительность при росте объема данных и не рассчитана на высокие нагрузки. PostgreSQL, напротив, является промышленным стандартом, обеспечивает полную поддержку транзакций, сложных агрегирующих запросов (необходимых для расчета суточной калорийности и статистики питания), надежное хранение данных и масштабируемость.
\end{itemize}

Для реализации клиентской части приложения выбран подход на основе HTML, CSS, Bootstrap и чистого JavaScript. При создании интерфейсов использовались официальные руководства из документации разработчиков~\cite{bootstrap_docs}.

Для реализации серверной части приложения выбран фреймворк Django (Python), архитектура которого подробно описана в работе~\cite{django_docs}.
Использование встроенной ORM позволяет эффективно обрабатывать запросы от множества пользователей при одновременном подсчете калорий и агрегации данных о питании.

\subsection{Формулировка цели и задач работы}

На основе проведенного анализа инструментов (см. подраздел 1.2) и выявленных
пробелов в существующих продуктах, была сформулирована цель исследования.

\textbf{Цель работы} --- создание веб-приложения для учета потребляемых продуктов, автоматического подсчета калорий и анализа структуры питания пользователей.


Для достижения поставленной цели необходимо решить следующие \textbf{задачи}:
\begin{enumerate}
	\item Провести сравнительный анализ существующих приложений для учета питания и методов расчета суточной нормы калорий;
	\item Разработать архитектуру клиент-серверного приложения и спроектировать структуру базы данных PostgreSQL для хранения информации о пользователях, продуктах и истории приемов пищи;
	\item Реализовать программные модули для ведения дневника питания (добавление/удаление записей) и мониторинга суточной калорийности;
	\item Интегрировать справочник продуктов и внедрить алгоритм автоматического подсчета КБЖУ на основе веса порции;
	\item Провести тестирование разработанной системы на корректность расчета калорий и удобство работы с дневником питания.
\end{enumerate}